\section{Maximal affine families of odd theta characteristics}
\label{sec:odd-theta-affine}

Let \(C\) be a smooth projective geometrically connected curve of genus
\(g\geq1\) over an algebraically closed field of characteristic different
from \(2\). A theta characteristic is a line bundle whose square is
the canonical bundle. Write \(J(C)\) for the Jacobian and
\(\Theta(C)\) for the theta characteristics of \(C\),
\(V=J(C)[2]\), and \(B:V\times V\to\mathbb F_2\) for the additive Weil
pairing. Thus \(\dim_{\mathbb F_2}V=2g\), \(B\) is nondegenerate
alternating, and \(\Theta(C)\) is a \(V\)-torsor. Set
\(\epsilon(\theta)=h^0(C,\theta)\bmod2\).
The parity theorem says
\begin{equation}
 \epsilon(\theta+u+v)+\epsilon(\theta+u)
 +\epsilon(\theta+v)+\epsilon(\theta)=B(u,v).
 \label{eq:theta-polar}
\end{equation}
The number of even and odd characteristics is respectively
\(2^{g-1}(2^g+1)\) and \(2^{g-1}(2^g-1)\).
For complex curves these are Atiyah's Theorems 2 and 3, printed
pp.~48--49 \citep{Atiyah1971Spin}; Mumford's second and third results,
printed p.~182, give the algebraic statement with characteristic
different from \(2\) \citep{Mumford1971Theta}. These are the only geometric
inputs needed below.

\paragraph{Quadratic refinements and coordinates.}
For \(\theta\in\Theta(C)\) put
\[
 q_\theta(v)=\epsilon(\theta+v)+\epsilon(\theta).
\]
Equation \eqref{eq:theta-polar} makes \(q_\theta\) a quadratic refinement
of \(B\), and
\begin{equation}
 q_{\theta+u}(v)=q_\theta(v)+B(u,v).
 \label{eq:theta-translation}
\end{equation}
Nondegeneracy of \(B\) shows that \(\theta\mapsto q_\theta\) bijects the
theta characteristics with all quadratic refinements of \(B\). For a
symplectic basis \(e_i,f_i\), writing \(v=\sum_i(a_ie_i+b_if_i)\) gives
\[
 q(v)=\sum_i a_ib_i+\sum_i a_iq(e_i)+\sum_i b_iq(f_i).
\]
The elementary two-variable character sum therefore gives
\begin{equation}
 \sum_{v\in V}(-1)^{q(v)}
   =2^g(-1)^{\sum_iq(e_i)q(f_i)}.
 \label{eq:theta-gauss}
\end{equation}
This identifies the basis-independent Arf invariant. The parity counts
above imply
\(\operatorname{Arf}(q_\theta)=\epsilon(\theta)\).
In particular, for every symplectic basis there is a unique
\(\theta_0\) with \(q_{\theta_0}(e_i)=q_{\theta_0}(f_i)=0\); it is even,
and
\begin{equation}
 \epsilon\!\left(\theta_0+\sum_i(a_ie_i+b_if_i)\right)=\sum_i a_ib_i.
 \label{eq:theta-hyperbolic-coordinates}
\end{equation}
All assertions below concern this finite affine quadratic space.
They make no assertion about geometric theta characteristics of curves
in characteristic \(2\).

\begin{theorem}[Classification and exact number]
\label{thm:maximal-odd-affine}
A nonempty affine family \(A=\theta+W\subseteq\Theta(C)\) consists of
odd theta characteristics precisely when \(\epsilon(\theta)=1\) and
\(q_\theta|_W=0\). Every inclusion-maximal such family has dimension
\(g-1\), and hence cardinality \(2^{g-1}\). Moreover:
\begin{enumerate}
\item Every \((g-1)\)-dimensional \(B\)-isotropic subspace \(W\subset V\)
is the direction of exactly one maximal odd affine family.
\item These families form a single orbit under the natural symplectic
action on quadratic refinements.
\item Their number is
\begin{equation}
 M_g=(2^g-1)\prod_{j=2}^g(2^j+1).
 \label{eq:odd-family-count}
\end{equation}
Each odd theta characteristic lies in exactly
\(\prod_{j=2}^g(2^j+1)\) such families.
\end{enumerate}
Consequently, for \(g=2k\) the maximum dimension is \(2k-1\), the
maximum cardinality is \(2^{2k-1}\), and the minimum ambient codimension
is \(2k+1\).
\end{theorem}

\begin{proof}
The first criterion follows directly from the definition of \(q_\theta\).
Its polarization implies \(B|_{W\times W}=0\), so \(r=\dim W\leq g\).
Extend a basis of a \(B\)-isotropic \(r\)-space \(W\) to a symplectic
basis, with \(W=\langle e_1,\ldots,e_r\rangle\), and use
\eqref{eq:theta-hyperbolic-coordinates}. An affine coset of \(W\) has
constant parity exactly when \(b_1=\cdots=b_r=0\).
Such cosets are indexed by the remaining \(2s\) coordinates, where
\(s=g-r\), and their parity is
\(\sum_{i=r+1}^g a_ib_i\). The character sum \eqref{eq:theta-gauss}
therefore counts their odd and even constant cosets as
\begin{equation}
 O(g,r)=2^{s-1}(2^s-1),\qquad E(g,r)=2^{s-1}(2^s+1)
 \quad(s\geq1).
 \label{eq:constant-coset-count}
\end{equation}
When \(s=0\), there is one constant even coset and no odd one.
Thus no odd family has dimension \(g\), and each \(W\) of dimension
\(g-1\) has exactly one odd coset.

It remains to prove inclusion-maximality, rather than only the largest
possible dimension. Suppose \(\theta+W\) is odd and \(r<g-1\).
Here \(q_\theta\) has Arf invariant \(1\) and vanishes on \(W\).
Choose a symplectic extension \(e_i,f_i\), \(1\leq i\leq r\), of a basis
of \(W\), and replace each \(f_i\) by
\(f_i+q_\theta(f_i)e_i\). This preserves the symplectic pairings and makes
\(q_\theta(f_i)=0\). Hence these \(r\) hyperbolic planes split off an
orthogonal complement of dimension \(2s=2(g-r)\) and Arf invariant \(1\).
It has
\[
 2^{2s-1}-2^{s-1}>1
\]
zeros, because \(s\geq2\), by \eqref{eq:theta-gauss}. A nonzero singular
vector \(v\) in that complement extends \(W\) to \(W+\langle v\rangle\)
while preserving \(q_\theta=0\). Therefore every lower-dimensional odd
affine family has a one-dimensional extension. The already proved
dimension bound makes all \((g-1)\)-dimensional odd families maximal.

For completeness, count \(B\)-isotropic \(r\)-spaces by ordered bases.
After choosing an isotropic \(i\)-space, its perpendicular contains
\(2^{2g-i}\) vectors; exclude its own \(2^i\) vectors to choose the next
basis vector. Dividing by the number of ordered bases of an \(r\)-space
gives
\begin{equation}
 I(g,r)=
 \prod_{i=0}^{r-1}
 \frac{2^{2g-i}-2^i}{2^r-2^i}
 =
 \frac{\prod_{j=g-r+1}^g(2^{2j}-1)}
      {\prod_{j=1}^r(2^j-1)}.
 \label{eq:isotropic-space-count}
\end{equation}
At \(r=g-1\) this simplifies to \eqref{eq:odd-family-count}.
The symplectic group is transitive on isotropic subspaces of fixed
dimension, by extension to symplectic bases. Its action
\(q\mapsto q\circ S^{-1}\) preserves the Arf invariant and transports
the unique odd coset attached to a direction; this proves the orbit
claim. This action is also transitive on Arf-invariant-one refinements:
successively split off hyperbolic planes using a nonzero singular vector
until only a two-dimensional anisotropic plane remains, which has a
unique nonzero-value quadratic form. Double-counting incident odd
points and maximal families gives
\[
 \frac{M_g\,2^{g-1}}{2^{g-1}(2^g-1)}
 =\prod_{j=2}^g(2^j+1)
\]
families through each odd point.
\end{proof}

\paragraph{Classical terminology and scope.}
An isotropic group of characteristics is a G\"opel group; its affine
cosets are G\"opel systems. The constant-coset counts
\eqref{eq:constant-coset-count}, including uniqueness of the odd coset
at rank \(g-1\), are explicitly recorded in
Bernatska--Kopeliovich, Section~2.6, p.~7
\citep{BernatskaKopeliovich2023}, which refers to Baker's Articles
297 and 299--300. Thus the classification is classical theta
combinatorics, and no novelty is claimed.

One can state exactly the relevant meaning of \emph{syzygetic}.
A triple \(\theta_1,\theta_2,\theta_3\) is syzygetic when
\[
 \epsilon(\theta_1)+\epsilon(\theta_2)+\epsilon(\theta_3)
 +\epsilon(\theta_1+\theta_2+\theta_3)=0;
\]
the sum of three points in an \(\mathbb F_2\)-affine torsor is intrinsic.
Every triple in an odd affine family is syzygetic. Conversely, if a
nonempty set \(S\) of odd characteristics has all its triples
syzygetic, fix \(\theta\in S\). The differences \(u,v\) from \(\theta\)
satisfy \(q_\theta(u)=q_\theta(v)=q_\theta(u+v)=0\), so \(B(u,v)=0\).
Their linear span is therefore \(q_\theta\)-singular, and the affine
span of \(S\) is entirely odd. Thus the maximal such syzygetic sets are
exactly the maximal affine families in the theorem. This statement
does not require theta characteristics to have exactly one section,
nor does it assert a configuration theorem for bitangent hyperplanes.

\begin{corollary}[Exact affine covering number]
\label{cor:odd-theta-cover}
Covering all odd theta characteristics by affine families contained
in the odd locus requires and suffices to use \(2^g-1\) families.
More generally, if such a cover contains at least a fraction \(\delta\)
of all odd characteristics, it has at least
\(\lceil\delta(2^g-1)\rceil\) members.
\end{corollary}
\begin{proof}
Each family has at most \(2^{g-1}\) points, whereas the odd locus has
\(2^{g-1}(2^g-1)\) points. This proves the lower bounds by the union
bound. For the upper bound, use \(q(a,b)=a\cdot b\) and, for each
nonzero \(b\in\mathbb F_2^g\), take
\[
 A_b=\{(a,b):a\cdot b=1\}.
\]
The \(2^g-1\) sets \(A_b\) are disjoint affine \((g-1)\)-flats and
partition the odd locus.
\end{proof}

\paragraph{Application to inner-product parity fibres.}
For \(k\geq1\) four-bit inner-product blocks put
\[
 P(y)=\sum_{i=1}^k(y_{i,0}y_{i,1}+y_{i,2}y_{i,3}).
\]
This is the hyperbolic quadratic form of genus \(g=2k\). If a
nonempty affine fibre \(F\subseteq\mathbb F_2^{4k}\) satisfies
\(P|_F=1\), then
\begin{equation}
 \dim F\leq2k-1,\qquad \operatorname{codim}F\geq2k+1.
 \label{eq:odd-ip-rank}
\end{equation}
Both inequalities are sharp. The even fibre has maximum affine
dimension \(2k\), so the bound there is only codimension \(2k\), also
sharp. In resolution over parity equations, denoted Res\((\oplus)\),
a line is a disjunction of affine equations over \(\mathbb F_2\).
Its rank is the rank of their coefficient matrix; its falsifying set
is the affine solution set obtained by negating every equation.
If \(F\) is obtained by fixing the outside coordinates of the
falsifying affine flat of a line of rank \(r\), its codimension is at
most \(r\); thus \(P|_F=1\) forces \(r\geq2k+1\).

This sharpens the isotropic rank estimate by one for an odd
wrong-parity fibre. A boundary condition \(P=c\) has wrong-parity fibre
\(P=1\) only when \(c=0\); when \(c=1\), the wrong-parity fibre is even
and the improvement by one does not apply. Moreover,
Corollary~\ref{cor:odd-theta-cover} is a cardinality statement with
respect to the uniform distribution on the full odd fibre. It supplies
neither a common boundary set for arbitrary proof lines nor a lower
bound on the mass of their union under a
distribution used in a proof-complexity argument. For a refutation $\pi$, a distribution $\mu$, falsifying sets $A_L$ of its lines, a rank threshold $\rho$, and a positive constant $\delta$, the required assertion would have the form
\begin{equation}
 \sum_{\substack{L\in\pi\\\operatorname{rank}L\geq\rho}}
       \mu(A_L)\geq\delta.
 \label{eq:theta-m3-not-proved}
\end{equation}
Establishing this covering inequality requires an additional argument. No
superpolynomial Res\((\oplus)\) lower bound follows without it.

\paragraph{Finite examples.}
The accompanying integer-bit program enumerates all symplectic
isotropic directions through genus \(4\), all odd cosets in each
direction, and every one-step extension of a lower-dimensional odd
flat. Through genus \(3\) it additionally enumerates all affine
directions, without imposing isotropy, and obtains the same odd
families. The resulting maximal-family counts are
\[
\begin{array}{c|rrrr}
g&1&2&3&4\\ \hline
\#\Theta_{\mathrm{odd}}&1&6&28&120\\
\text{maximum affine dimension}&0&1&2&3\\
\#\text{maximal families}&1&15&315&11475 .
\end{array}
\]
These finite checks illustrate the preceding proof; the proof applies in every genus.

